\begin{textsample}{NEG Dim 1 – al – Score 6.00 – al/al\_ef\_97.txt}  \label{ex:f1_neg_140}
5.3 As Unidades Temáticas A partir da BNCC ( 2017 ), e com o Referencial Curricular de Alagoas ( 2014 ), a área de Ciências da Natureza se concentra em estudar os fenômenos naturais, sendo assim, envolve os processos de evolução e manutenção da vida no planeta, perpassando – com seus recursos naturais, transformações e fontes de energias – do nosso planeta no Sistema Solar e no Universo, bem como a aplicabilidade do conhecimento científico nas várias esferas da vida humana ( BRASIL, 2017 ).
Todo esse conhecimento se materializa como desenvolvimento de estudantes capazes de compreender, se comunicar e intervir nos processos de sistematização de conceitos, ampliando a visão de mundo.
A estruturação da organização curricular da área, componente de Ciências, as aprendizagens essenciais se encontram em três unidades temáticas, que se repetem ao longo do Ensino Fundamental.
Desta forma, as Unidades Temáticas são : Matéria e energia ; Vida e evolução ; Terra e universo.
Cada Unidade Temática compõe -se de Objetos de Conhecimentos que são os conteúdos conceituais, assim como traz consigo, também, as habilidades pertinentes à unidade e ao objeto que podem ser organizados pelo professor no seu planejamento, atendendo às especificidades socioculturais locais.
Essa organização indica perspectivas de abordagem para trabalhar a especificidade da área de forma progressiva, sem isolar, pois estabelece diferentes sequências nos ciclos, permite tratar as temáticas de importância local e fazer conexão entre os objetos de conhecimento da área.
Esses devem ser vistos não de forma horizontal apenas, mas na verticalidade, e no processo de ensino de forma espiral, em que conhecimentos se entrelacem num todo para a formação integral do estudante.

% matched lemmas: 
\end{textsample}
